\section{Síntesis de los principales hallazgos}

La investigación permitió implementar un prototipo funcional que integra procesamiento de datos crediticios, entrenamiento automatizado, registro de experimentos, almacenamiento analítico y visualización en el contexto de una institución de la economía social y solidaria.\\

Los resultados obtenidos son específicos al conjunto de datos, la configuración experimental y la división temporal utilizada, por lo que no deben generalizarse a otras instituciones o entornos sin una validación adicional.\\ 

En la evaluación presentada en este estudio, LightGBM mostró una mayor capacidad discriminativa que las implementaciones de CNN y MLP bajo las configuraciones documentadas. El rendimiento de LightGBM disminuyó conforme aumentó el horizonte, y el recall presentó una reducción especialmente relevante en los plazos largos.\\

El análisis del código fuente permitió documentar ampliamente las ocho condiciones de \texttt{crisis\_flag}, las 21 características, la arquitectura exacta de cada modelo y el protocolo de división de datos (49\% entrenamiento, 21\% validación, 30\% prueba).

\section{Respuesta al objetivo general}

El objetivo general se alcanzó en lo relativo al diseño y la integración técnica de los componentes principales. Se implementaron flujos de entrenamiento e inferencia en Apache Airflow, registro en MLflow, almacenamiento en PostgreSQL y dashboards en Apache Superset. El resultado debe presentarse como un prototipo funcional en el entorno evaluado. La evidencia disponible no permite concluir que la plataforma controle la mora, reduzca pérdidas o haya sido validada para apoyar decisiones en un entorno productivo; dicha validación operativa y la evaluación de impacto en las decisiones quedan a cargo de la institución objeto del presente trabajo. El resultado debe presentarse como un prototipo funcional. La evidencia disponible no permite concluir que la plataforma controle la mora, reduzca pérdidas o haya sido validada para apoyar decisiones en un entorno productivo.

\section{Respuesta a los objetivos específicos}

Respecto al primer objetivo, se identificaron factores plausibles asociados con el comportamiento de las redes neuronales: una muestra secuencial reducida (solo bloques con $\geq 24$ meses de historial generan secuencias), desbalance de clases sin compensación en MLP, ventanas de seis meses y configuración documentada completa. El MLP no implementa \texttt{sample\_weights}, lo que puede explicar su tendencia a predecir todas las observaciones como ``no crisis''. No se realizaron experimentos de ablación, por lo que estos factores no pueden declararse causas demostradas.\\

En relación con el segundo objetivo, LightGBM alcanzó el mayor AUC-ROC bajo las configuraciones disponibles dentro del contexto de ejecución de este estudio. Esta comparación no permite afirmar que los modelos basados en árboles sean superiores en todos los conjuntos tabulares, ya que la búsqueda de hiperparámetros no fue equivalente entre modelos (LightGBM utiliza semilla fija y configuración documentada; CNN y MLP no fijan semilla), y los resultados dependen del conjunto de datos y la configuración experimental.\\

Para el tercer objetivo, se observó una disminución del rendimiento de LightGBM con el horizonte. El AUC-ROC pasó de 0,874 en el horizonte de 1 mes a 0,702 en el de 18 meses, mientras que el recall descendió de 0,542 a 0,034. El análisis detallado (Tabla A.1) revela que, a partir del horizonte $h = 8 $, el recall cae a 0, lo que significa que el modelo no identifica ningún caso de crisis en horizontes intermedios y largos bajo el umbral de 0,5. Por tanto, aunque el sistema fue diseñado para predecir 18 horizontes, los resultados indican que su utilidad operativa se concentra en los horizontes cortos ($ h \le 6 $). Los horizontes largos deben interpretarse con cautela y no pueden considerarse operativamente útiles sin calibración, selección de umbral y análisis de costos de error.\\

El cuarto objetivo se cumplió como implementación técnica. El datamart integra dimensiones de tiempo, riesgo, sector y sucursal, y los dashboards permiten consultar información histórica y predicciones. No se evaluó la utilidad con usuarios, la latencia ni el impacto sobre decisiones, por lo que esas dimensiones permanecen pendientes.

\section{Contribuciones}

La primera contribución es una arquitectura de extremo a extremo que enlaza ETL, modelado, MLOps, datamart y visualización, demostrando la viabilidad técnica de integrar estos componentes en el contexto de una institución de la economía social y solidaria. La segunda es la comparación de tres enfoques para dieciocho horizontes, junto con el análisis de la degradación temporal. La tercera es la implementación de un esquema estrella y dashboards que integran predicciones y agregados históricos.\\

Las contribuciones se circunscriben al prototipo y a la institución analizada. No se afirma superioridad universal de LightGBM ni se generalizan los resultados a otras entidades financieras.

\section{Aspectos éticos}

El desarrollo del presente trabajo involucra datos crediticios reales de una institución financiera. Por ello, se consideraron los siguientes aspectos éticos:

\begin{itemize}
    \item \textbf{Confidencialidad}: Los datos fueron anonimizados a nivel de bloque-mes, sin incluir identificadores personales. El acceso al repositorio está restringido al autor y al tutor.
    \item \textbf{Transparencia}: El modelo seleccionado (LightGBM) permite calcular importancias de características, lo que facilita la interpretación de las predicciones. No obstante, no se realizó una auditoría formal de sesgos por sector económico o región.
    \item \textbf{Impacto en decisiones}: El prototipo no reemplaza el juicio de los analistas de crédito; su función es complementaria y su uso debe estar sujeto a validación humana. La institución debe establecer políticas claras sobre el uso de predicciones automáticas antes de cualquier implementación productiva.
\end{itemize}

\section{Limitaciones}

El estudio presenta limitaciones metodológicas y de generalización. La definición de \texttt{crisis\_flag} quedó documentada mediante ocho condiciones con pesos acumulativos (umbral $\geq 5$), pero existe un riesgo confirmado de circularidad de constructo (ver Sección 4.5): las variables judiciales (\texttt{tasa\_judicial}, \texttt{creditos\_judiciales}, \texttt{total\_costo\_judicial}) se utilizan tanto para definir la etiqueta como características de entrada. Esta dependencia puede inflar artificialmente el poder predictivo de esas variables y del modelo.\\

La división de datos sigue un protocolo temporal con tres conjuntos (49\% entrenamiento, 21\% validación, 30\% prueba), y el early stopping utiliza exclusivamente el conjunto de validación. Sin embargo, no se reportan conteos por horizonte, no se fijan semillas para CNN y MLP, y no se realizan repeticiones para estimar variabilidad.\\

No se reportan PR-AUC, calibración, curvas por umbral ni análisis detallado de falsos negativos. La comparación entre modelos utiliza presupuestos de ajuste no equivalentes. Además, el estudio se realizó en una sola institución y no incluye validación externa.\\

Desde la perspectiva operativa, las capturas demuestran la implementación de los flujos y dashboards, pero no se presentan pruebas de rendimiento, disponibilidad, recuperación ante fallos, integridad de datos o usabilidad. Estas limitaciones impiden considerar el prototipo como un sistema productivo validado.

\section{Trabajo futuro}

El trabajo futuro se organiza en cuatro líneas prioritarias:

\begin{enumerate}
    \item \textbf{Reproducibilidad y robustez}: Registrar conteos por etapa y horizonte, fijar semillas para todos los modelos, realizar repeticiones con múltiples semillas (al menos 10), y versionar datos, código y entornos mediante \texttt{dvc} o \texttt{git-lfs}.
    
    \item \textbf{Análisis de circularidad de constructo}: Repetir la evaluación excluyendo las variables judiciales (\texttt{tasa\_judicial}, \texttt{creditos\_judiciales}, \texttt{total\_costo\_judicial}) para estimar cuánto del poder predictivo se debe a la dependencia con la etiqueta. Este es el experimento más crítico antes de cualquier implementación operativa.
    
    \item \textbf{Comparación equilibrada}: Equilibrar los presupuestos de ajuste (añadir \texttt{sample\_weights} al MLP, realizar búsqueda de hiperparámetros equivalente para CNN y MLP), reportar intervalos de confianza, incluir PR-AUC, calibración y selección de umbrales basada en costos de error (falsos positivos vs. falsos negativos).
    
    \item \textbf{Validación externa}: Evaluar el prototipo con datos de al menos otra institución financiera y mediante pruebas funcionales con usuarios analistas. Solo después de estas validaciones, que incluyan la evaluación de usabilidad, tiempos de respuesta e impacto en los procesos de decisión, sería posible considerar el sistema como una herramienta operativa para la gestión de cartera.
    
\end{enumerate}